\documentclass[12pt]{article}
\usepackage[utf8]{inputenc}
\usepackage{amsmath}
\usepackage{geometry}
\usepackage{booktabs}
\usepackage{tabularx}
 \geometry{letterpaper, margin=0.2in}

\begin{document}

\section*{The Governing Heat Equation}
For a one-dimensional rod, the heat conduction is governed by:
\[
\rho c_p \frac{\partial T}{\partial t} = k \frac{\partial^2 T}{\partial x^2}
\]
\begin{itemize}
    \item $\frac{\partial T}{\partial t} = $ the rate of change of temperature with respect to time ($^\circ C / s$)
    \item $\frac{\partial^2 T}{\partial x^2} = $ sensitivity of temperature changes at a specific point from its neighbors ($^\circ C / m^2$)
    \item $t =$ Time (s), $\quad \bullet$ $x = $ Position along the rod $(m)$, $\quad \bullet$ $\alpha = \frac{k}{\rho c_p}$ is the thermal diffusivity $(m^2/s)$
    \item $\rho =$ Density $(kg/m^3)$ $\quad \bullet$ $k =$ Thermal Conductivity $(W/m\cdot^\circ C)$, $\quad \bullet$ $c_p =$ Specific Heat $(J/kg\cdot^\circ C)$
\end{itemize}

\section*{Boundary and Initial Conditions}
The model uses the Dirichlet boundary conditions
\begin{flalign*}
    & \hspace{1.5cm} T(0, t) = 100^\circ\text{F} \quad \text{(Left Boundary)} & \\
    & \quad\quad\quad T(L, t) = 20^\circ\text{F} \quad \text{(Right Boundary)} & \\
    & \quad\quad\quad T(x, 0) = T(x) \quad \text{(Initial Temperature)} &
\end{flalign*}

\section*{Finite Element Analysis}
The rod is divided into $n$ elements of length $h = L/n$. \\
sFor each element, the local stiffness matrix $\mathbf{K}^{(e)}$ \\
and force vector $\mathbf{F}^{(e)}$ are defined as:
\begin{flalign*}
&\textbf{1. Spatial Discretization (FEM):} & \\
&\text{The rod is divided into } n \text{ elements of length } h. \\ 
&\text{ The local connectivity is defined by:} & \\
&\mathbf{K}^{(e)} = \frac{k}{h} \begin{bmatrix} 1 & -1 \\ -1 & 1 \end{bmatrix} & \\
&\textbf{2. Transient System (Backward Euler):} & \\
&\text{To ensure unconditional stability, an implicit scheme is used:} & \\
&(\mathbf{I} - \alpha \Delta t \mathbf{K}') \mathbf{T}^{n+1} = \mathbf{T}^{n} + \mathbf{b}_{bc} & \\
&\textbf{3. Boundary Treatment (Dirichlet):} & \\
&\text{The vector } \mathbf{b}_{bc} \text{ enforces the fixed temperatures at the boundaries } (100^\circ\text{F}, 20^\circ\text{F}), & \\
&\text{effectively driving the heat flow toward a linear steady-state where } \mathbf{K}\mathbf{T} = \mathbf{0}. &
\end{flalign*}

\vspace{10cm}
Accounting for ambient air influence, the 1D heat conduction is governed by:
\[
\frac{\partial T}{\partial t} = \alpha \frac{\partial^2 T}{\partial x^2} - \nu(T - T_\infty)
\]
\begin{itemize}
    \item $\nu$: Convection Coefficient $(s^{-1})$ $\quad \bullet$ $T =$ Temperature ($^\circ$C), $\quad \bullet$ $T_{\infty} =$ Ambient Temperature ($^\circ$C)
\end{itemize}

\[
\frac{\partial T}{\partial t} = \alpha \frac{\partial^2 T}{\partial x^2} - \nu(T - T_\infty)
\quad
\begin{aligned}
    &\bullet \nu: \text{Convection Coefficient } (s^{-1}) \\ 
    &\bullet T = \text{Temperature } (^{\circ}\text{C}) \\ 
    &\bullet T_{\infty} = \text{Ambient Temperature } (^{\circ}\text{C})
\end{aligned}
\]

\vspace{10cm}


\begin{table}[h]
\centering
\label{tab:material_comparison}
\begin{tabularx}{\textwidth}{lXXX}
\toprule
\textbf{Property} & \textbf{Copper (Bulk)} & \textbf{Additive (3D-Printed)} & \textbf{Physical Interpretation} \\
\midrule
Conductivity ($k$) & $384.72$ W/m$\cdot$K & $14.40$ W/m$\cdot$K & XN is $\approx$ 26x more resistive to heat flow than Cu. \\
\addlinespace
Diffusivity ($\alpha$) & $0.0738$ m$^2$/s & $0.0213$ m$^2$/s & Cu reaches steady state $\approx$ 3.5x faster than XN. \\
\addlinespace
Convection ($h$) & $165.41$ W/m$^2\cdot$K & $128.70$ W/m$^2\cdot$K & Variations due to ambient lab/wind conditions. \\
\addlinespace
Avg. Std. Error & $0.95$ & $0.73$ & Low error validates 1D FEA model robustness. \\
\bottomrule
\end{tabularx}
\end{table}




\begin{table}[h]
\centering
\caption{Thermal Parameter Comparison: Bulk Copper vs. 3D-Printed Steel (XN)}
\label{tab:material_comparison}
% We use 'l' for the labels and 'c' for the numbers to keep them tight. 
% The final 'X' column will take up all the remaining horizontal space.
\begin{tabularx}{\textwidth}{l c c X}
\toprule
\textbf{Property} & \textbf{Copper} & \textbf{3D Printed} & \textbf{Physical Interpretation} \\
\midrule
Conductivity ($k$) & $384.72$ W/m$\cdot^{\circ}\text{C}$ & $14.40$ W/m$\cdot^{\circ}\text{C}$ & 3DP is $\approx$ 26x more resistive to heat flow than Cu. \\
\addlinespace
Diffusivity ($\alpha$) & $0.0738$ m$^2$/s & $0.0213$ m$^2$/s & Cu reaches steady state $\approx$ 3.5x faster than 3DP. \\
\addlinespace
Convection ($h$) & $165.41$ W/m$^2\cdot^{\circ}\text{C} $& $128.70$ W/m$^2\cdot^{\circ}\text{C}$ & Variations due to ambient lab/wind conditions. \\
\addlinespace
Avg. Std. Error & $0.95$ & $0.73$ & Low error validates 1D FEA model robustness. \\
\bottomrule
\end{tabularx}
\end{table}

\end{document}